\newpage
\pagenumbering{arabic}
\setcounter{page}{1}

\section{Вовед}

Во современите работни простории се посветува сè поголемо внимание врз
квалитетот на работната околина и здравствените навики на работниците,
особено кај инженерите и ИТ професионалците кои се долго време во
затворени простории како и пред компјутер на своите бироа. Најразличните
параметри како температура, влажност на воздухот, светлина,
концентрација на различни гасови како CO2, NO2, CO, бучавата и
осветлувањето во просторијата имаат значајна улога во продуктивноста,
фокусот па и општата благосостојба на човекот, но во најголемиот дел од
компаниите тие не се следат континуирано или објективно. Токму оваа
потреба, мене лично, а и на другите околу мене за подобро и подетално
разбирање на условите во работната средина претставува основна
мотивација за развој на оваа платформа која овозможува автоматизирано
прибирање, визуелизација и анализа на релевантни податоци.

На пазарот постојат повеќе решенија за IoT мониторинг, но секое од нив
има свои ограничувања. Home Assistant е популарна open-source платформа
која овозможува интеграција на различни паметни уреди, но бара
комплексна конфигурација и техничко знаење за поставување, а не
обезбедува вградени механизми за безбедна автентикација на уредите.
Комерцијалните решенија како AWS IoT Core и Azure IoT Hub нудат напредни
enterprise функционалности, решенија и висока скалабилност, но нивната
цена и комплексност ги прават непрактични за индивидуална употреба или
мали организации. ThingsBoard претставува добра платформа за
визуелизација на IoT податоци, но нема вградена интеграција со модели за
вештачка интелигенција за напредна анализа. Од друга страна,
здравствените платформи како Google Fit, Samsung Health, Mi Fitness и
други, работат во затворени екосистеми и не дозволуваат едноставна
интеграција со околински сензори или сопствени системи.

Иако постојат многу поединечни решенија за мониторинг на одделни
параметри, или пак можност за интеграција на разни уреди, ниту еден од
нив не претставува интегриран систем кој истовремено ги спојува
околинските услови, здравствените навики на корисниците како и
надворешните услови како временска прогноза или квалитетот на воздухот.
Дополнително, се соочуваат со разни проблеми како недоволно интуитивен и
напреден кориснички интерфејс, небезбеден начин на испраќање на
информациите, како и неефикасен начин на складирање на собраните
податоци. Овие недостатоци создаваат простор и потреба за посеопфатно и
безбедно решение.

Целта на овој дипломски труд е да се изработи скалабилна, безбедна и
интуитивна платформа за автоматизирано прибирање, складирање,
визуелизација и анализа на податоци од сите достапни извори, со што се
добива увид со широк спектар за работната околина, условите во неа, како
и за самиот работник. Специфичните цели на проектот се:

\begin{itemize}
\item
  Дизајн и имплементација на микросервисна архитектура која овозможува
  модуларност и независно скалирање на компонентите
\item
  Имплементација на безбедна комуникација помеѓу уредите и платформата
  со взаемна TLS (mTLS) автентикација и X.509 сертификати
\item
  Развој на ефикасен систем за прием и складирање на временски податоци
  користејќи \texttt{Redis Streams} и \texttt{TimescaleDB}
\item
  Интеграција на сервис за анализа и интерпретација на собраните
  податоци со помош вештачка интелигенција и машинско учење
\item
  Изработка на интуитивна веб контролна табла со прилагодливи widgets за
  визуелизација и приказ на информациите
\item
  Развој на Android апликација за прибирање и испраќање на здравствени и
  фитнес податоци
\end{itemize}

Во втората глава од овој дипломски труд ја разгледувам теоретската
основа потребна за разбирање на системот: IoT концептите, MQTT
протоколот, безбедносните механизми со mTLS и X.509 сертификати,
микросервисната архитектура и временските бази на податоци. Во третата
глава го претставувам дизајнот на системот, вклучувајќи ја
архитектурата, протокот на податоци, безбедносниот модел и дизајнот на
базата на податоци. Четвртата глава е посветена на имплементацијата на
сите компоненти: сервисите за управување со уреди, прием на MQTT пораки,
запишување во база, вештачка интелигенција, веб контролната табла и
Android апликацијата. Во петтата глава ги прикажувам резултатите од
тестирањето и работата на системот, а во шестата глава го сумирам
заклучокот и предлагам идни насоки за развој.
